% 原书物理页=印刷页35--43上（第3章章前言与 3.1 全节；3.1.6 内命题的证明跨至 p43 上部，止于 3.2 节标题前）.
\chapter{半单 Lie 代数的结构}\label{chap:ssl-3}
\enlargethispage{1pt}% avoid a sub-point output-box overflow on the chapter opening
在本部分中，$K$ 表示一个特征为零的代数闭域.

\section{Cartan 子代数与根}\label{sec:ssl-3-1}
设 $\mathfrak g$ 是一个 $K$-Lie 代数. 若 $\mathfrak h\subset\mathfrak g$ 是一个 Lie 子代数，记 $N(\mathfrak h):=\{X\in\mathfrak g,\ \operatorname{ad}_X(\mathfrak h)\subset\mathfrak h\}$ 为它的规范化子，它也是 $\mathfrak g$ 的 Lie 子代数（请验证）.
\begin{definition*}
$\mathfrak g$ 的一个\newterm{ssl-cartan-subalgebra}{Cartan 子代数}{Cartan subalgebra}是指 $\mathfrak g$ 的一个幂零子代数 $\mathfrak h$，使得 $N(\mathfrak h)=\mathfrak h$.
\end{definition*}
我们来证明这样的子代数确实存在. 对 $X\in\mathfrak g$，记 $\mathfrak g_{X,0}:=\bigcup_{n\in\mathbb N}\operatorname{Ker}(\operatorname{ad}_X)^n\subset\mathfrak g$，即 $\operatorname{ad}_X$ 的属于特征值 $0$ 的广义特征空间. 注意 $\dim_K(\mathfrak g_{X,0})$ 是 $0$ 在 $\operatorname{ad}_X$ 的特征多项式中的重数. 若 $\mathfrak g_{X,0}$ 按包含关系在诸 $\mathfrak g_{Z,0}$（$Z\in\mathfrak g$）中极小，就称 $X$ 为\newterm{ssl-quasi-regular}{拟正则元素}{quasi-regular element}.
\begin{proposition*}
若 $X$ 拟正则，则 $\mathfrak g_{X,0}$ 是 $\mathfrak g$ 的 Cartan 子代数. 反之，$\mathfrak g$ 的任何 Cartan 子代数都具有这种形状.
\end{proposition*}
\begin{proof}
i) 注意 $\mathfrak g_{X,0}$ 总是一个 Lie 子代数. 这确实由下述（可用归纳法证明的）公式得出：
\[(\operatorname{ad}_X)^n([Y,Z])=\sum_{k=0}^n\binom{n}{k}\big[(\operatorname{ad}_X)^k(Y),(\operatorname{ad}_X)^{n-k}(Z)\big].\]
此外，$X\in\mathfrak g_{X,0}$，故 $\operatorname{ad}_X(N(\mathfrak g_{X,0}))\subset\mathfrak g_{X,0}$. 由此推出 $(\operatorname{ad}_X)^{\dim(\mathfrak g_{X,0})+1}(N(\mathfrak g_{X,0}))=0$，于是 $N(\mathfrak g_{X,0})\subset\mathfrak g_{X,0}$，最终 $N(\mathfrak g_{X,0})=\mathfrak g_{X,0}$.

ii) 现在设 $X$ 拟正则，证明 $\mathfrak g_{X,0}$ 是幂零的. 只须证明对每个 $Y\in\mathfrak g_{X,0}$，自同态 $(\operatorname{ad}_Y)|_{\mathfrak g_{X,0}}$ 都是幂零的. 为简化记号，置 $\mathfrak h:=\mathfrak g_{X,0}$、$r:=\dim(\mathfrak g_{X,0})$. 我们把 $\operatorname{ad}_Y$ 在 $\operatorname{End}_K(\mathfrak g)$ 中的特征多项式记作 $\Phi^{\mathfrak g}_{\operatorname{ad}Y}\in K[T]$. 由于 $\operatorname{ad}_Y$ 稳定 $\mathfrak h$，有因式分解 $\Phi^{\mathfrak g}_{\operatorname{ad}Y}=\Phi^{\mathfrak h}_{\operatorname{ad}Y}\,\Phi^{\mathfrak g/\mathfrak h}_{\operatorname{ad}Y}$，我们要证 $\Phi^{\mathfrak h}_{\operatorname{ad}Y}=T^r$. 取定 $\mathfrak h$ 的一组基 $X_1,\cdots,X_r$，并写 $Y=\sum_i y_iX_i$，$y_i\in K$. 那么，由于 $K$ 无穷，存在唯一的诸多项式 $\Phi_{\mathfrak g},\Phi_{\mathfrak h}$ 与 $\Phi_{\mathfrak g/\mathfrak h}\in K[Y_1,\cdots,Y_r,T]$，它们在特化 $Y_i\mapsto y_i$ 下给出 $\Phi^{\mathfrak g}_{\operatorname{ad}Y},\Phi^{\mathfrak h}_{\operatorname{ad}Y}$ 与 $\Phi^{\mathfrak g/\mathfrak h}_{\operatorname{ad}Y}$. 在唯一分解整环 $K[Y_1,\cdots,Y_r,T]$ 中也有因式分解 $\Phi_{\mathfrak g}=\Phi_{\mathfrak h}\Phi_{\mathfrak g/\mathfrak h}$. 我们必须证明 $\Phi_{\mathfrak h}=T^r$. 记 $\operatorname{ord}_T(f)$ 为多项式 $f$ 的 $T$-赋值（即使得 $T^k$ 整除 $f$ 的最大整数 $k$）. 由于 $\Phi_{\mathfrak h}$ 是关于 $T$ 的、系数在 $K[Y_1,\cdots,Y_r]$ 中的 $r$ 次首一多项式，只须证明 $\operatorname{ord}_T(\Phi_{\mathfrak h})=r$. 注意，按 $\mathfrak h$ 的定义，在 $K[T]$ 中 $T$ 不整除 $\Phi^{\mathfrak g/\mathfrak h}_{\operatorname{ad}X}$. 由于 $\Phi^{\mathfrak g/\mathfrak h}_{\operatorname{ad}X}$ 是 $\Phi_{\mathfrak g/\mathfrak h}$ 的一个特化，由此推出在 $K[Y_1,\cdots,Y_r,T]$ 中 $T$ 不整除 $\Phi_{\mathfrak g/\mathfrak h}$，因而 $\operatorname{ord}_T(\Phi_{\mathfrak h})=\operatorname{ord}_T(\Phi_{\mathfrak g})$. 由于 $K$ 无穷，可找到特化 $Y_i\mapsto z_i$，使得（记 $Z=\sum_i z_iX_i$）$\operatorname{ord}_T(\Phi^{\mathfrak h}_{\operatorname{ad}Z})=\operatorname{ord}_T(\Phi_{\mathfrak h})$ 且 $\operatorname{ord}_T(\Phi^{\mathfrak g}_{\operatorname{ad}Z})=\operatorname{ord}_T(\Phi_{\mathfrak g})$. 于是等式 $\operatorname{ord}_T(\Phi^{\mathfrak h}_{\operatorname{ad}Z})=\operatorname{ord}_T(\Phi^{\mathfrak g}_{\operatorname{ad}Z})$ 蕴涵 $\mathfrak g_{Z,0}\subset\mathfrak h=\mathfrak g_{X,0}$. 但由于 $X$ 是正则的\footnote{译注：原文此处作``$X$ est régulier''（正则）；本命题假设的是 $X$ 拟正则，且正则元素迟至下一个定义才引入，疑为``拟正则''之误.}，由此得 $\mathfrak g_{Z,0}=\mathfrak h=\mathfrak g_{X,0}$，进而 $\operatorname{ord}_T(\Phi^{\mathfrak g}_{\operatorname{ad}Z})=r$，最终得到所要的 $\operatorname{ord}_T(\Phi_{\mathfrak h})=r$.

iii) 现在设 $\mathfrak h$ 是 $\mathfrak g$ 的一个 Cartan 子代数. 取定 $\mathfrak h$ 的一组基 $X_1,\cdots,X_r$，并如上考虑 $K[Y_1,\cdots,Y_r,T]$ 中的因式分解 $\Phi_{\mathfrak g}=\Phi_{\mathfrak h}\Phi_{\mathfrak g/\mathfrak h}$. 由于 $\mathfrak h$ 幂零，我们知道 $\Phi_{\mathfrak h}=T^r$. 此外，若 $T$ 整除 $\Phi_{\mathfrak g/\mathfrak h}$，则 $0$ 将是 $\mathfrak h$ 在 $\mathfrak g/\mathfrak h$ 中的一个权，于是由 Engel 定理，$\mathfrak g/\mathfrak h$ 将有一个被 $\mathfrak h$ 零化的非零元素 $\bar Z$. 这意味着，若 $Z\in\mathfrak g$ 是 $\bar Z$ 的一个原像，则 $Z\notin\mathfrak h$ 且 $[\mathfrak h,Z]\subset\mathfrak h$. 但这样一来 $Z$ 就规范化了 $\mathfrak h$，与假设 $\mathfrak h=N(\mathfrak h)$ 矛盾. 故 $T$ 不整除 $\Phi_{\mathfrak g/\mathfrak h}$. 适当特化后，得到一个元素 $X\in\mathfrak h$，使 $T$ 不整除 $\Phi^{\mathfrak g/\mathfrak h}_{\operatorname{ad}X}$. 这蕴涵 $\mathfrak g_{X,0}\subset\mathfrak h$. 由于 $\mathfrak h$ 幂零，又有 $\mathfrak h\subset\mathfrak g_{X,0}$，最终 $\mathfrak h=\mathfrak g_{X,0}$.
\end{proof}
Cartan 子代数的存在性既已获得，我们关心的是它在共轭意义下的唯一性. 注意，甚至先验地都不清楚所有 Cartan 子代数是否具有相同的维数.
\begin{definition*}
称整数 $\operatorname{rk}(\mathfrak g):=\min_{X\in\mathfrak g}\big(\dim_K(\mathfrak g_{X,0})\big)$ 为 $\mathfrak g$ 的\newterm{ssl-rank}{秩}{rank}. 若 $\dim(\mathfrak g_{X,0})=\operatorname{rk}(\mathfrak g)$，则称元素 $X\in\mathfrak g$ 为\newterm{ssl-regular-element}{正则元素}{regular element}.
\end{definition*}
于是每个正则元素都是拟正则的，且相应的 Cartan 子代数 $\mathfrak g_{X,0}$ 的维数是 $r=\operatorname{rk}(\mathfrak g)$.
\begin{exercise*}
设 $\mathfrak g=\mathfrak{gl}_n(K)$，$\mathfrak b$（相应地 $\mathfrak b_{\mathrm{nilp}}$）是上三角矩阵（相应地幂零上三角矩阵）构成的子代数.
\begin{enumerate}[label=\roman*)]
\item 利用 $[\mathfrak b,\mathfrak b]\subset\mathfrak b_{\mathrm{nilp}}$，证明 $\operatorname{rk}(\mathfrak g)\geqslant n$.
\item 若 $X=\sum_{i=1}^n x_iE_{ii}$ 且诸 $x_i$ 两两不同，证明 $\mathfrak g_{X,0}$ 是对角矩阵构成的子代数. 由此得出 $\operatorname{rk}(\mathfrak g)=n$. 再证明 $\operatorname{rk}(\mathfrak{sl}_n(K))=n-1$.
\item 证明 $X\in\mathfrak{gl}_n(K)$ 正则当且仅当它的特征值两两不同.
\end{enumerate}
\end{exercise*}
\begin{lemma*}
$\mathfrak g$ 的正则元素所成的集合 $\mathfrak g_{\mathrm{reg}}$ 是 $\mathfrak g$ 的一个非空\newterm{ssl-zariski-open}{Zariski 开集}{Zariski open set}. 特别地，当 $K=\mathbb C$ 时，它对解析拓扑（即有限维 $\mathbb C$-向量空间的拓扑）而言是 $\mathfrak g$ 的稠密连通开集.
\end{lemma*}
\begin{proof}
如前一个证明那样选取 $\mathfrak g$ 的一组基，并把``特征多项式''函数 $X\mapsto\Phi_{\operatorname{ad}X}$ 看作一个元素 $\Phi\in K[Y_1,\cdots,Y_{\dim\mathfrak g},T]$. 那么 $r$ 就是 $\Phi$ 的 $T$-赋值，从而存在非零多项式 $a_r=a_r(Y_1,\cdots,Y_{\dim\mathfrak g})\in K[Y_1,\cdots,Y_{\dim\mathfrak g}]$，使得
\[\Phi\equiv a_rT^r\ \text{modulo}\ (T^{r+1}).\]
于是元素 $X$ 正则当且仅当 $a_r(X)\neq 0$，其中 $a_r(X)$ 表示多项式 $a_r$ 在 $X$ 于所选基下的坐标处的取值. 这样，$\mathfrak g_{\mathrm{reg}}$ 是一个 Zariski 开集，并且由于 $K$ 代数闭（Hilbert 零点定理），它非空. 当 $K=\mathbb C$ 时，与仿射空间的任何 Zariski 开集一样，它对解析拓扑是稠密且道路连通的开集.
\end{proof}
\begin{theorem*}
设 $K=\mathbb C$，并记 $\operatorname{Int}(\mathfrak g)$ 为 $\operatorname{Aut}\mathfrak g$ 中由诸 $\exp(\operatorname{ad}_X)$（$X\in\mathfrak g$）生成的闭子群. 则 $\mathfrak g$ 的所有 Cartan 子代数在 $\operatorname{Int}(\mathfrak g)$ 下互相共轭.
\end{theorem*}
\begin{remark*}
把 $\operatorname{Int}(\mathfrak g)$ 换成由诸 $\exp(\operatorname{ad}_X)$（其中 $\operatorname{ad}_X$ 幂零，从而使指数确有定义）生成的子群后，本定理对任何特征为零的代数闭域 $K$ 都成立. 那时的证明是初等的，但冗长而艰难. 在 $K=\mathbb C$ 的情形，我们将给出一个相当短的证明，不过要用到一点微分几何.
\end{remark*}
\begin{remark*}
本定理也蕴涵每个拟正则元素都是正则的.
\end{remark*}
\begin{proof}
关键一点如下. 设 $\mathfrak h$ 是一个 Cartan 子代数，$\mathfrak h_{\mathrm{qreg}}$ 是它的拟正则元素所成的集合；由前一命题证明的末尾（它是使 $T$ 不整除 $\Phi^{\mathfrak g/\mathfrak h}_{\operatorname{ad}X}$ 的点集），这是一个 Zariski 开集（从而是解析开集）. 简单地记 $G:=\operatorname{Int}(\mathfrak g)$；由于是闭的，它是 $\mathrm{GL}(\mathfrak g)$ 的一个 Lie 子群. 那么，映射
\[G\times\mathfrak h_{\mathrm{qreg}}\longrightarrow\mathfrak g,\quad (g,X)\mapsto g\cdot X\]
是\newterm{ssl-submersion}{淹没}{submersion}的.

暂且承认这一点. 第一个推论是：这个映射的像是开集，从而与稠密开集 $\mathfrak g_{\mathrm{reg}}$ 的交非平凡. 由此 $\mathfrak h_{\mathrm{reg}}:=\mathfrak h\cap\mathfrak g_{\mathrm{reg}}$ 非空，且由于是 Zariski 开集，它是 $\mathfrak h$ 的稠密开集. 于是考虑 $\mathfrak g_{\mathrm{reg}}$ 上由下式定义的等价关系
\[X\sim Y\quad\text{若}\ \mathfrak g_{X,0}\ \text{与}\ \mathfrak g_{Y,0}\ \text{共轭}.\]
对 $X\in\mathfrak h_{\mathrm{reg}}$，上述淹没性表明，存在 $X$ 在 $\mathfrak g_{\mathrm{reg}}$ 中的一个由与 $X$ 等价的元素组成的开邻域. 这样，上述等价关系的诸等价类都是开集. 但由于 $\mathfrak g_{\mathrm{reg}}$ 连通，等价类只有一个，定理随之得出.

剩下要证上述淹没性. 借助 $G$ 中的平移，只须证明切映射在诸点 $(e,X)$（$X\in\mathfrak h_{\mathrm{qreg}}$）处的满射性. 按定义，$G$ 在 $e$ 处的切空间 $\operatorname{Lie}(G)$ 是 $\mathfrak{gl}(\mathfrak g)$ 中由诸 $\operatorname{ad}_Y$（$Y\in\mathfrak g$）生成的 Lie 子代数. 于是 $\operatorname{Lie}(G)=\operatorname{ad}(\mathfrak g)$. 同样按定义，切映射由下式给出
\[f\colon \operatorname{ad}(\mathfrak g)\times\mathfrak h\longrightarrow\mathfrak g,\quad (\operatorname{ad}_Y,H)\mapsto\operatorname{ad}_Y(X)+H=-\operatorname{ad}_X(Y)+H.\]
它的像因此包含 $\mathfrak h$. 此外，按 $\mathfrak h_{\mathrm{qreg}}$ 的定义，$\operatorname{ad}_X$ 在 $\mathfrak g/\mathfrak h$ 上的作用可逆，故 $f$ 的像可满射到 $\mathfrak g/\mathfrak h$ 上. 由此 $f$ 是满射.
\end{proof}
\begin{example*}
在 $\mathfrak{gl}_n(\mathbb C)$ 中，任何 Cartan 子代数都在 $\mathrm{GL}_n(\mathbb C)$ 下共轭于由可对角化矩阵组成的子代数\footnote{译注：原文作``matrices diagonalisables''（可对角化矩阵）；可对角化矩阵一般不构成子代数，此处应指由对角矩阵（matrices diagonales）组成的子代数.}.
\end{example*}
\begin{exercise*}
对 $\mathfrak{sl}_n$ 与 $\mathfrak{sp}_{2n}$ 证明类似的命题，并由此推出 $\operatorname{rk}(\mathfrak{sp}_{2n})=n$. 对 $\mathfrak{so}_n$，对角矩阵构成的子代数是零！但可以更换非退化对称双线性型：代替标准形式，在偶数维考虑形式 $\psi_{2n}(v,w)=2\sum_{i=1}^n v_iw_{i+n}$，在奇数维考虑形式 $\psi_{2n+1}(v,w)=2\sum_{i=1}^n v_iw_{i+n}+v_{2n+1}w_{2n+1}$. 那么 $\mathfrak{so}(\psi_m)$ 同构于 $\mathfrak{so}_m$，并且可以证明对角矩阵构成的子代数是一个 Cartan 子代数. 于是可以看到 $\operatorname{rk}(\mathfrak{so}_{2n})=\operatorname{rk}(\mathfrak{so}_{2n+1})=n$.
\end{exercise*}

\begin{numpara}[Cartan 子代数的根]\label{np:ssl-3-1-1}
设 $\mathfrak g$ 是一个 $K$-Lie 代数，$\mathfrak h\subset\mathfrak g$ 是一个 Cartan 子代数. 让 $\mathfrak h$ 通过 $\mathfrak g$ 的伴随表示 $(\mathfrak g,\operatorname{ad})$ 在 $\mathfrak h$ 上的限制作用在 $\mathfrak g$ 上. 对每个线性型 $\lambda\in\mathfrak h^{*}_{\mathrm{ab}}$，有广义特征空间
\[\mathfrak g^\lambda=\{Y\in\mathfrak g,\ \forall X\in\mathfrak h,\ (\operatorname{ad}_X-\lambda(X))^n(Y)=0\ \text{对}\ n\geqslant\dim\mathfrak g\ \text{成立}\}.\]
\begin{lemma*}
有 $\mathfrak g^0=\mathfrak h$，且对任何 $\lambda,\mu\in\mathfrak h^{*}_{\mathrm{ab}}$，有 $[\mathfrak g^\lambda,\mathfrak g^\mu]\subset\mathfrak g^{\lambda+\mu}$.
\end{lemma*}
\begin{proof}
由于 $\mathfrak h$ 幂零，对每个 $X\in\mathfrak h$，$(\operatorname{ad}_X)|_{\mathfrak h}$ 幂零，故 $\mathfrak h\subset\mathfrak g^0$. 反之，$\mathfrak h$ 以幂零算子作用在 $\mathfrak g^0/\mathfrak h$ 上. 于是若 $\mathfrak g^0\supsetneq\mathfrak h$，就存在 $X\in\mathfrak g^0\setminus\mathfrak h$，它在 $\mathfrak g^0/\mathfrak h$ 中的像 $\bar X$ 被 $\mathfrak h$ 零化. 但这就意味着 $[\mathfrak h,X]\subset\mathfrak h$，即 $X\in N(\mathfrak h)=\mathfrak h$：矛盾.

第二个断言由公式
\[(\operatorname{ad}_X-\lambda(X)-\mu(X))^n([Y,Z])=\sum_{k=0}^n\binom{n}{k}\big[(\operatorname{ad}_X-\lambda(X))^k(Y),(\operatorname{ad}_X-\mu(X))^{n-k}(Z)\big]\]
得出，留给读者.
\end{proof}
\begin{definition*}
$\mathfrak h$ 在 $\mathfrak g$ 中的一个\newterm{ssl-root}{根}{root}是指 $\mathfrak h$ 在 $(\mathfrak g,\operatorname{ad})$ 中的一个非零权. 我们把根的全体记作 $\Phi$.
\end{definition*}
关于幂零代数的表示的定理给我们一个\newterm{ssl-root-decomposition}{根分解}{root decomposition}
\[\mathfrak g=\bigoplus_{\alpha\in\Phi\cup\{0\}}\mathfrak g^\alpha,\]
而前面的命题告诉我们：对 $\alpha,\beta\in\Phi$，括积 $[\mathfrak g^\alpha,\mathfrak g^\beta]$ 在 $\alpha+\beta\notin\Phi\cup\{0\}$ 时为零，在 $\alpha+\beta=0$ 时含于 $\mathfrak h$，在 $\alpha+\beta\in\Phi$ 时含于 $\mathfrak g^{\alpha+\beta}$.
\begin{example*}
考察 $\mathfrak g=\mathfrak{sl}_n(K)$ 而 $\mathfrak h$ 为对角矩阵构成的 Cartan 子代数的情形. 公式
\[[\sum_i x_{ii}E_{ii},\ \sum_{i,j}a_{ij}E_{ij}]=\sum_{i,j}(x_{ii}-x_{jj})a_{i,j}E_{ij}\]
表明 $\Phi=\{\alpha_{ij}\in\mathfrak h^*,\ 1\leqslant i\neq j\leqslant n\}$，其中
\[\alpha_{ij}\Big(\sum_i x_{ii}E_{ii}\Big)=x_{ii}-x_{jj},\]
并且 $\mathfrak g^{\alpha_{ij}}=\mathfrak g_{\alpha_{ij}}=KE_{ij}$. 于是 $\mathfrak h$ 的广义权空间都是名副其实的权空间，且都是一维的. 我们将看到，对半单代数而言这是一个普遍事实.
\end{example*}

\textbf{从现在起假设 $\mathfrak g$ 是半单的.}

我们记 $\langle X,Y\rangle=B_{\operatorname{ad}}(X,Y)$ 为 $\mathfrak g$ 的 Killing 型，它因此是非退化的.
\end{numpara}
\begin{proposition}\label{prop:ssl-3-1-2}
设 $\mathfrak g$ 半单，$\mathfrak h$ 是它的一个 Cartan 子代数. 那么：
\begin{enumerate}[label=\roman*)]
\item 对 $\Phi\cup\{0\}$ 中满足 $\alpha+\beta\neq 0$ 的 $\alpha,\beta$，有 $\langle\mathfrak g^\alpha,\mathfrak g^\beta\rangle=0$；
\item $\langle.,.\rangle|_{\mathfrak h}$ 非退化；
\item $\langle.,.\rangle|_{\mathfrak g^\alpha\times\mathfrak g^{-\alpha}}$ 诱导 $\mathfrak g^\alpha$ 与 $\mathfrak g^{-\alpha}$ 之间的一个对偶.
\end{enumerate}
\end{proposition}
\begin{proof}
i) 设 $\lambda\in\Phi\cup\{0\}$、$X\in\mathfrak g^\alpha$、$Y\in\mathfrak g^\beta$. 那么 $\operatorname{ad}_X\operatorname{ad}_Y(\mathfrak g^\lambda)\subset\mathfrak g^{\lambda+\alpha+\beta}$，而 $\lambda+\alpha+\beta\neq\lambda$. 由此，在从属于分解 $\mathfrak g=\bigoplus_\lambda\mathfrak g^\lambda$ 的一组基下，$\operatorname{ad}_X\operatorname{ad}_Y$ 的矩阵对角线为零，故 $\operatorname{tr}(\operatorname{ad}_X\operatorname{ad}_Y)=0$.

ii) 设 $X\in\mathfrak h$. 由 i)，对每个根 $\alpha\in\Phi$ 有 $X\in\mathfrak g^{\alpha,\perp}$. 于是若 $X\in\mathfrak h^\perp$，则 $X\in\mathfrak g^\perp$，从而 $X=0$.

iii) 设 $X\in\mathfrak g^\alpha$. 由 i)，对每个权 $\beta\neq-\alpha$ 有 $X\in\mathfrak g^{\beta,\perp}$. 于是若 $X\in\mathfrak g^{\perp-\alpha}$，则 $X\in\mathfrak g^\perp$，从而 $X=0$.
\end{proof}
\begin{remark*}
$\langle.,.\rangle$ 在 $\mathfrak h$ 上限制的非退化性与 $\mathfrak h$ 的幂零性并不矛盾. 事实上，由于 $\mathfrak h$ 不是理想，这一限制先验地（事后也一样）与 $\mathfrak h$ 自身的 Killing 型毫无关系.
\end{remark*}
\begin{exercise*}
证明：若把 Killing 型换成任何一个非退化双线性型\footnote{译注：原文仅设``非退化''（non-d\'eg\'en\'er\'ee）；i) 的证明还需该双线性型是不变的（invariante），疑脱``不变''一词.}，命题的结论仍然成立.
\end{exercise*}
\begin{corollary*}
假设同上. 那么 $\mathfrak h$ 是 Abel 的，由半单元素组成，且 $\Phi$ 张成对偶空间 $\mathfrak h^*$.
\end{corollary*}
\begin{proof}
i) 取 $\mathfrak g$ 的一个旗 $\mathcal V$，使 $\operatorname{ad}(\mathfrak h)\subset\mathfrak b(\mathcal V)$（由于 $\mathfrak h$ 可解）. 于是 $\operatorname{ad}([\mathfrak h,\mathfrak h])\subset\mathfrak b^{\mathrm{nilp}}(\mathcal V)$. 设 $Y\in[\mathfrak h,\mathfrak h]$. 对每个 $X\in\mathfrak h$，有 $\operatorname{ad}_X\operatorname{ad}_Y\in\mathfrak b^{\mathrm{nilp}}(\mathcal V)$，故 $\operatorname{tr}(\operatorname{ad}_X\operatorname{ad}_Y)=0$. 于是 $Y\in\mathfrak h^\perp$，由前一命题的 ii) 得 $Y=0$. 故 $[\mathfrak h,\mathfrak h]=0$，$\mathfrak h$ 确是 Abel 的.

ii) 设 $X\in\mathfrak h$. 要证 $X=X_s$，即 $X_n=0$. 由于 $\mathfrak h$ 等于它的规范化子，它在 Jordan 分解下稳定，故 $X_n\in\mathfrak h$. 但由于 $\mathfrak h$ 是 Abel 的，对每个 $Y\in\mathfrak h$，$\operatorname{ad}_Y\operatorname{ad}_{X_n}$ 幂零，从而迹为零. 换句话说，$X_n\in\mathfrak h^\perp$，故 $X_n=0$.

iii) 设 $Y$ 满足对每个 $\alpha\in\Phi$ 有 $\alpha(Y)=0$. 那么 $\operatorname{ad}_Y$ 的特征值全为零，故 $\operatorname{ad}_Y$ 幂零. 如上可证 $Y\in\mathfrak h^\perp$，而命题的 ii) 给出 $Y=0$.
\end{proof}
\begin{exercise*}
证明当 $\mathfrak g$ 约化时本推论仍然成立，除了``$\Phi$ 张成 $\mathfrak h^*$''这一点.
\end{exercise*}
推论：由于每个 $\operatorname{ad}_X$（$X\in\mathfrak h$）都可对角化且 $\mathfrak h$ 是 Abel 的，$\mathfrak h$ 在 $\mathfrak g$ 上的伴随作用可以同时对角化. 于是对每个根 $\alpha$，有 $\mathfrak g_\alpha=\mathfrak g^\alpha$.
\begin{proposition}\label{prop:ssl-3-1-3}
对每个 $\alpha\in\Phi$，有 $\dim(\mathfrak g_\alpha)=1$ 且 $\{n\in\mathbb Z,\ n\alpha\in\Phi\}=\{\pm 1\}$.
\end{proposition}
\begin{proof}
设 $E_\alpha\in\mathfrak g_\alpha$. 于是对每个 $H\in\mathfrak h$ 有 $[H,E_\alpha]=\alpha(H)E_\alpha$. 取 $F\in\mathfrak g_{-\alpha}$ 使 $\langle E_\alpha,F\rangle=1$，置 $H'_\alpha:=[E_\alpha,F]\in\mathfrak h=\mathfrak g^0$. 那么对每个 $H\in\mathfrak h$ 有
\begin{equation}\tag{$(*)$}
\langle H,H'_\alpha\rangle=\langle H,[E_\alpha,F]\rangle=\langle[H,E_\alpha],F\rangle=\alpha(H)\langle E_\alpha,F\rangle=\alpha(H).
\end{equation}
证明 $\alpha(H'_\alpha)\neq 0$. 事实上，上面推论证明的 iii) 保证存在 $\beta$ 使 $\beta(H'_\alpha)\neq 0$. 于是考虑 $W:=\sum_{n\in\mathbb Z}\mathfrak g_{\beta+n\alpha}$. 那么 $W$ 在 $\operatorname{ad}_{E_\alpha}$ 与 $\operatorname{ad}_F$ 下稳定，从而也在 $\operatorname{ad}_{[E_\alpha,F]}=\operatorname{ad}_{H'_\alpha}$ 下稳定，并且 $\operatorname{tr}\big((\operatorname{ad}_{H'_\alpha})|_W\big)=\operatorname{tr}\big([(\operatorname{ad}_{E_\alpha})|_W,(\operatorname{ad}_F)|_W]\big)=0$. 而
\begin{align*}
\operatorname{tr}\big((\operatorname{ad}_{H'_\alpha})|_W\big)
&=\sum_{n\in\mathbb Z}\operatorname{tr}\big((\operatorname{ad}_{H'_\alpha})|_{\mathfrak g_{\beta+n\alpha}}\big)\\
&=\sum_{n\in\mathbb Z}(\beta+n\alpha)(H'_\alpha)\,\dim(\mathfrak g_{\beta+n\alpha}).
\end{align*}
由于 $\beta(H'_\alpha)\neq 0$，故必有 $\alpha(H'_\alpha)\neq 0$.

证明 $\operatorname{ad}_{E_\alpha}(\mathfrak g_{-\alpha})\subset K.H'_\alpha$. 事实上，对每个 $Y\in\mathfrak g_{-\alpha}$ 与每个 $H\in\mathfrak h$，有 $\langle H,[E_\alpha,Y]\rangle=\langle[H,E_\alpha],Y\rangle=\alpha(H)\langle E_\alpha,Y\rangle$，故由 $(*)$ 得 $[E_\alpha,Y]=\langle E_\alpha,Y\rangle H'_\alpha$.

现在置
\[V:=K.E_\alpha\oplus K.H'_\alpha\oplus\Big(\bigoplus_{n<0}\mathfrak g_{n\alpha}\Big).\]
按构造，$V$ 在 $\operatorname{ad}_F$ 下稳定，且我们刚看到它也在 $\operatorname{ad}_{E_\alpha}$ 下稳定. 如上所述，由此推出它在 $\operatorname{ad}_{H'_\alpha}$ 下稳定，且 $\operatorname{tr}\big((\operatorname{ad}_{H'_\alpha})|_V\big)=0$. 而 $\operatorname{tr}\big((\operatorname{ad}_{H'_\alpha})|_V\big)=\alpha(H'_\alpha)\big(1+\sum_{n<0}n\,\dim(\mathfrak g_{n\alpha})\big)$. 由于 $\alpha(H'_\alpha)\neq 0$，由此便推出定理.
\end{proof}
\begin{corollary*}
对 $H,H'\in\mathfrak h$，有
\[\langle H,H'\rangle=\sum_{\alpha\in\Phi}\alpha(H)\alpha(H').\]
\end{corollary*}
\begin{proof}
显然，因为 $\operatorname{ad}_H\operatorname{ad}_{H'}(\mathfrak g_\alpha)\subset\mathfrak g_\alpha$，且 $\mathfrak g_\alpha=\mathfrak g_\alpha$\footnote{译注：原文两处均作下标 $\mathfrak g_\alpha$；按上文（见分目~\ref{np:ssl-3-1-1} 的例及本页``$\mathfrak g_\alpha=\mathfrak g^\alpha$''）第二处应为上标 $\mathfrak g^\alpha$，意为该（广义）权空间是一维权空间.} 是一维的.
\end{proof}
\begin{corollary}\label{cor:ssl-3-1-4}
$K$-向量空间 $\mathfrak{sl}_\alpha:=\mathfrak g_{-\alpha}\oplus[\mathfrak g_\alpha,\mathfrak g_{-\alpha}]\oplus\mathfrak g_\alpha$ 是 $\mathfrak g$ 的一个同构于 $\mathfrak{sl}_2(K)$ 的 Lie 子代数. 更精确地说：
\begin{itemize}
\item 存在唯一的 $H_\alpha\in[\mathfrak g_\alpha,\mathfrak g_{-\alpha}]$，使得 $\alpha(H_\alpha)=2$；
\item 一旦取定 $E_\alpha\in\mathfrak g_\alpha\setminus\{0\}$，就存在唯一的 $F_\alpha\in\mathfrak g_{-\alpha}$，使得 $[E_\alpha,F_\alpha]=H_\alpha$.
\end{itemize}
此时三元组 $(E_\alpha,H_\alpha,F_\alpha)$ 是一个\newterm{ssl-sl2-triplet}{$\mathfrak{sl}_2$ 三元组}{sl2-triplet}.
\end{corollary}
\begin{proof}
沿用前面证明中的记号，有 $\operatorname{ad}_{E_\alpha}(\mathfrak g_{-\alpha})=K.H'_\alpha$，故 $[\mathfrak g_\alpha,\mathfrak g_{-\alpha}]=K.H'_\alpha$ 是一维的. 因此存在唯一的 $H_\alpha\in[\mathfrak g_\alpha,\mathfrak g_{-\alpha}]$ 使 $\alpha(H_\alpha)=2$. 具体写出，就是 $H_\alpha=\frac{2H'_\alpha}{\alpha(H'_\alpha)}=\frac{2H'_\alpha}{\langle H'_\alpha,H'_\alpha\rangle}$. 现在存在唯一的 $F_\alpha\in\mathfrak g_{-\alpha}$ 使 $\langle E_\alpha,F_\alpha\rangle=\frac{2}{\langle H'_\alpha,H'_\alpha\rangle}$. 于是 $[E_\alpha,F_\alpha]=H_\alpha$，三元组 $(E_\alpha,H_\alpha,F_\alpha)$ 是一个 $\mathfrak{sl}_2$ 三元组.
\end{proof}
\begin{example*}
在 $\mathfrak{sl}_n$ 的情形，沿用前面例子中的记号，对任何 $i\neq j$ 有 $[\mathfrak g_{\alpha_{ij}},\mathfrak g_{\alpha_{ji}}]=K(E_{ii}-E_{jj})$，且 $(E_{ij},E_{ii}-E_{jj},E_{ji})$ 是一个 $\mathfrak{sl}_2$ 三元组.
\end{example*}
\begin{proposition}\label{prop:ssl-3-1-5}
设 $\beta\in\Phi$ 满足 $\beta\neq\pm\alpha$. 那么
\begin{enumerate}[label=\roman*)]
\item $\beta(H_\alpha)\in\mathbb Z$，且 $\beta-\beta(H_\alpha)\in\Phi$\footnote{译注：原文此处作 $\beta-\beta(H_\alpha)$；由下文证明及分目~\ref{np:ssl-3-1-6} iii) 应为 $\beta-\beta(H_\alpha)\alpha$ 之误.}. 此外，$\{n\in\mathbb Z,\ \beta+n\alpha\in\Phi\}$ 是 $\mathbb Z$ 中的一个区间；
\item 若 $\beta+\alpha\in\Phi$，则 $[\mathfrak g_\alpha,\mathfrak g_\beta]=\mathfrak g_{\alpha+\beta}$.
\end{enumerate}
\end{proposition}
\begin{proof}
考虑 $\mathfrak g_{\beta+\mathbb Z\alpha}:=\bigoplus_{n\in\mathbb Z}\mathfrak g_{\beta+n\alpha}$. 它是 $\mathfrak g$ 的一个在 $\operatorname{ad}_{E_\alpha}$ 与 $\operatorname{ad}_{F_\alpha}$ 下稳定的子空间，因而是 $\mathfrak{sl}_\alpha$ 的一个表示. $H_\alpha$ 在 $\mathfrak g_{\beta+\mathbb Z\alpha}$ 中的权是 $\{\beta(H_\alpha)+2n,\ n\ \text{使得}\ \mathfrak g_{\alpha+n\beta}\neq 0\}$\footnote{译注：原文此处下标作 $\mathfrak g_{\alpha+n\beta}$；按上式 $\mathfrak g_{\beta+\mathbb Z\alpha}$ 的定义应为 $\mathfrak g_{\beta+n\alpha}$.}.

$\mathfrak{sl}_2(K)$ 的表示理论告诉我们这些权都是整数，故 $\beta(H_\alpha)\in\mathbb Z$. 它还告诉我们：由于这些权的重数都是 $1$ 且奇偶性相同，$\mathfrak g_{\beta+\mathbb Z\alpha}$ 是 $\mathfrak{sl}_\alpha$ 的一个不可约表示，其权也形如 $\{-m,-m+2,\cdots,m-2,m\}$. 由此得到 i) 的\newterm{ssl-interval-property}{区间性质}{interval property}，并且由于 $\beta(H_\alpha)$ 是一个权，$-\beta(H_\alpha)$ 也是一个. 故 $\mathfrak g_{\beta-\beta(H_\alpha)\alpha}\neq 0$，即 $\beta-\beta(H_\alpha)\alpha\in\Phi$.

ii) 设 $\mathfrak g_{\beta+\alpha}\neq 0$. 它就是 $H_\alpha$ 的属于权 $\beta(H_\alpha)+2$ 的权空间. 由 $\mathfrak{sl}_2$ 不可约表示的已知结构，有 $[\mathfrak g_\alpha,\mathfrak g_\beta]=\operatorname{ad}_{E_\alpha}(\mathfrak g_\beta)=\mathfrak g_{\alpha+\beta}$，因为 $\mathfrak g_\beta$ 是 $H_\alpha$ 的属于权 $\beta(H_\alpha)$ 的权空间.
\end{proof}
\begin{corollary*}
对 $\alpha\in\Phi$，有 $\Phi\cap K\alpha=\{\pm\alpha\}$.
\end{corollary*}
\begin{proof}
假设存在 $\beta\in\Phi\cap K\alpha$ 异于 $\alpha$ 与 $-\alpha$. 由于 $\beta(H_\alpha)\in\mathbb Z$，有 $\beta\in\frac{1}{2}\mathbb Z\alpha$. 写 $\beta=\frac{n}{2}\alpha$，$n$ 为奇数. 那么由于 $-\beta\in\Phi$，命题~\ref{prop:ssl-3-1-5} i) 的区间性质保证 $\frac{1}{2}\alpha\in\Phi$. 但由命题 2.5.5\footnote{译注：原文作 2.5.5；本册 2.5 节仅有 2.5.1--4，此编号不存在. 本处所用根的整数倍性质见命题~\ref{prop:ssl-3-1-3}.}，$\Phi$ 中 $\frac{1}{2}\alpha$ 的倍数只有 $\pm\frac{1}{2}\alpha$，于是得到矛盾，因为 $\alpha\in\Phi$.
\end{proof}

\begin{numpara}[$\Phi$ 的有理性性质]\label{np:ssl-3-1-6}
总结一下到目前为止我们对集合 $\Phi$ 证明了些什么：
\begin{enumerate}[label=\roman*)]
\item $\Phi$ 张成 $K$-向量空间 $\mathfrak h^*$；
\item 对每个根 $\alpha\in\Phi$，有 $\Phi\cap K\alpha=\{\pm\alpha\}$；
\item 对两个根 $\beta,\alpha$，有 $\beta(H_\alpha)\in\mathbb Z$，且当 $\beta\neq-\alpha$ 时 $\{n\in\mathbb Z,\ \beta+n\alpha\in\Phi\}$ 是一个含 $-\beta(H_\alpha)$ 的区间.
\end{enumerate}
现在置
\[
\mathfrak h_{\mathbb Q}:=\operatorname{Vect}_{\mathbb Q}\{H_\alpha,\ \alpha\in\Phi\},
\]
以及
\[
\mathfrak h^{*}_{\mathbb Q}:=\operatorname{Vect}_{\mathbb Q}\{\alpha,\ \alpha\in\Phi\}.
\]
\begin{proposition*}
采用上述记号，则：
\begin{enumerate}[label=\roman*)]
\item $\dim_{\mathbb Q}(\mathfrak h_{\mathbb Q})=\dim_{\mathbb Q}(\mathfrak h^{*}_{\mathbb Q})=\dim_K(\mathfrak h)$；
\item 配对 $\mathfrak h\times\mathfrak h^{*}\longrightarrow K$ 诱导配对 $\mathfrak h_{\mathbb Q}\times\mathfrak h^{*}_{\mathbb Q}\longrightarrow\mathbb Q$；
\item Killing 型诱导一个有理形式 $\langle.,.\rangle\colon\mathfrak h_{\mathbb Q}\times\mathfrak h_{\mathbb Q}\longrightarrow\mathbb Q$；
\item $\langle.,.\rangle$ 到 $\mathfrak h_{\mathbb R}$ 的扩张是正定的.
\end{enumerate}
\end{proposition*}
\begin{proof}
i) 取 $\alpha_1,\cdots,\alpha_r$ 为 $\mathfrak h^{*}$ 的一组由 $\Phi$ 中元素组成的基. 那么诸 $H_{\alpha_i}$ 是 $\mathfrak h$ 的一组 $K$-基（注意 $\langle H_{\alpha_i},H\rangle=\frac{\langle H_{\alpha_i},H_{\alpha_i}\rangle}{2}\alpha_i(H)=\frac{2}{\langle H'_{\alpha_i},H'_{\alpha_i}\rangle}\alpha_i(H)$）. 对 $\alpha\in\Phi$，写 $H_\alpha=\sum_{i=1}^r\lambda_iH_{\alpha_i}$，$\lambda_i\in K$. 我们要证明对一切 $i$ 有 $\lambda_i\in\mathbb Q$. 设 $M\in M_r(K)$ 为矩阵 $(\langle H_{\alpha_i},H_{\alpha_j}\rangle)_{1\leqslant i,j\leqslant r}$. 那么 $M$ 可逆（因为 $\langle.,.\rangle$ 非退化），而 $\lambda=(\lambda_1,\cdots,\lambda_r)\in K^r$ 是线性方程组 $M.\lambda=\mu_\alpha$ 的唯一解，其中 $\mu_\alpha=(\langle H_{\alpha_1},H_\alpha\rangle,\cdots,\langle H_{\alpha_r},H_\alpha\rangle)\in K^r$. 而对两个根 $\alpha,\beta$ 有
\[\langle H_\alpha,H_\beta\rangle=\sum_{\gamma\in\Phi}\gamma(H_\alpha)\gamma(H_\beta)\in\mathbb Z,\]
故 $M\in M_r(\mathbb Q)$、$\mu_\alpha\in\mathbb Q^r$，最终 $\lambda\in\mathbb Q^r$. 由此推出诸 $H_{\alpha_i}$（$i=1,\cdots,r$）构成 $\mathfrak h_{\mathbb Q}$ 的一组 $\mathbb Q$-基. 同理可证 $\alpha_1,\cdots,\alpha_r$ 构成 $\mathfrak h^{*}_{\mathbb Q}$ 的一组 $\mathbb Q$-基.

ii) $\mathfrak h$ 与 $\mathfrak h^{*}$ 之间的 $K$-对偶确实把 $\mathfrak h_{\mathbb Q}\times\mathfrak h^{*}_{\mathbb Q}$ 映入 $\mathbb Q$，因为对一切 $\alpha,\beta\in\Phi$ 有 $\beta(H_\alpha)\in\mathbb Z$. 由于 $\mathfrak h_{\mathbb Q}$（相应地 $\mathfrak h^{*}_{\mathbb Q}$）在 $K$ 上张成 $\mathfrak h$（相应地 $\mathfrak h^{*}$），该配对是完美的.

iii) 同理. 我们已看到 $\langle\mathfrak h_{\mathbb Q},\mathfrak h_{\mathbb Q}\rangle\subset\mathbb Q$，且由于 $\mathfrak h_{\mathbb Q}$ 张成 $\mathfrak h$，该配对非退化.

iv) 由于 $\langle H,H\rangle=\sum_\alpha\alpha(H)^2$，所考虑的实双线性型是半正定的. 而已知它非退化，故它还是正定的.
\end{proof}
\end{numpara}
